% Copyright 2026 The terCAD team. All rights reserved.
% Use of this content is governed by a CC BY-NC-ND 4.0 license that can be found in the LICENSE file.

Dear readers, \\
I invite you to join me on a journey into mobile application development, starting from the
ground up. This means that we may not know anything about mobile development, native toolchains
(Android SDK, Xcode), the \q{NativePHP} framework, or how the \q{PHP} language, which usually
lives on a server, could be packed into a pocket device.

This book offers a different perspective on concepts you are interested in, inspired by, or
practicing, based on the author's vision. I invite you to share your thoughts and help me
identify any errors or typos in the repository that contains the book's sources and the
application we will develop together as we progress through the book. You can find the
repository at \href{https://github.com/lyskouski/app-game-fantasyland}{https://github.com/lyskouski/app-game-fantasyland}.
And, who knows?.. Maybe your insights will enrich the core ideas of this book, and we will
publish a new edition together!

Returning to the book's core themes and selected technologies, it's clear that businesses want
to reuse their existing expertise and codebase to reach users on their phones while avoiding the
cost of maintaining separate teams for each platform. For years, PHP developers (and the
\q{Laravel} community in particular) had to transition to Swift, Kotlin, or JavaScript-based
frameworks to deliver a mobile experience. Recently, NativePHP has emerged as a bridge that
has been used to embed a PHP runtime directly into iOS and Android applications, allowing them
to be built with the tools that are already known. This book, "\emph{From Zero to Market with
NativePHP: Mobile App Development and Distribution}," leads us through the process of creating a
real application: a mobile client for the long-running browser game Liga of Heroes
(\href{https://fantasyland.ru}{fantasyland.ru}). The application has been published on the Apple
App Store, Google Play, Samsung Galaxy Store, and Huawei AppGallery.

The main aim of this book is to demonstrate how to take advantage of the features of NativePHP
to create a mobile application with many functions on top of an existing web-based service.
Throughout the book, we will explore every phase of the application development process, from
conceptualization to distribution:


\begin{itemize}
\setlength{\itemsep}{3pt}
\setlength{\parskip}{0pt}
\setlength{\parsep}{0pt}
  \item Take basics in \q{PHP}-language, \q{Laravel}-framework and \q{NativePHP}, and how they
  turn a web-application into a native one.

  \item Prepare the development environment (PHP, Composer, Node.js, Android SDK, Xcode) and
  resolve the issues that appear on the first startup.

  \item Gain insights into the design principles of a client application that relies on an
  external service.

  \item Go through User Interface and Experience practices to adapt a desktop-oriented game for
  a small touch screen.

  \item Discover how to incorporate essential features into the application, and extend it by
  native plugins (device screen, storage, notifications).

  \item Explore strategies for secure interaction with external systems, session handling, and
  ensuring data integrity.

  \item Master techniques for testing and debugging practices, and automate them by Continuous
  Integration.

  \item Configure the application distribution across mobile platforms (Android, iOS) to their
  marketplaces, and explore strategies for ongoing maintenance and updates.
\end{itemize}

\noindent So, let's embark on this learning curve journey and let our PHP go native.
